\section*{Capítulo \ref{ch:b3:renal-osmoregulation} --- Fisiologia renal e osmorregulação}

\begin{solution}{exo:b3:renal-osmoregulation:1}
\omterm{def:b3:renal-osmoregulation:nephron}{Glomérulo}: filtra o plasma menos as proteínas. \omterm{def:b3:renal-osmoregulation:nephron}{Túbulo proximal}:
reabsorve dois terços do sal e da água, toda a glicose e os aminoácidos,
o bicarbonato. Ramo descendente de Henle: perde água para a medula. Ramo
ascendente: bombeia NaCl para fora, impermeável à água. Túbulo distal:
ajusta o sódio (\omterm{def:b3:renal-osmoregulation:controls}{aldosterona}), secreta potássio. \omterm{def:b3:renal-osmoregulation:nephron}{Ducto coletor}: reabsorção
de água sob a vasopressina, manejo de prótons e de ureia.
\end{solution}

\begin{solution}{exo:b3:renal-osmoregulation:2}
Clearance: o volume de plasma depurado da substância por unidade de
tempo, $UV/P$. Aqui, $100\times 1.5/2 = \qty{75}{mL/min}$, abaixo da TFG
de \qty{120}{mL/min}: a substância é filtrada e em parte reabsorvida
(cerca de \qty{38}{\%} da quantidade filtrada).
\end{solution}

\begin{solution}{exo:b3:renal-osmoregulation:3}
Um peixe de água doce é mais salgado que seu meio: a água entra pelas
brânquias por osmose, de modo que beber só aumentaria a inundação; ele
excreta urina copiosa e diluída, e as \omterm{def:b3:renal-osmoregulation:comparative}{células de cloreto} de suas
brânquias captam Na$^{+}$ e Cl$^{-}$ ativamente da água. Um peixe marinho
é menos salgado que o mar: perde água pelas brânquias, bebe água do mar
para repô-la, e as \omterm{def:b3:renal-osmoregulation:comparative}{células de cloreto} bombeiam para fora o sal engolido.
\end{solution}

\begin{solution}{exo:b3:renal-osmoregulation:4}
O \omterm{prop:b3:renal-osmoregulation:countercurrent}{efeito único} é a diferença de \qty{200}{mOsm/L} que o ramo ascendente
espesso mantém entre o interstício e seu próprio líquido, bombeando NaCl
para fora sem deixar que a água o siga. O ramo ascendente precisa ser
impermeável à água, ou a água seguiria o sal e a diferença
desapareceria; o ramo descendente precisa ser permeável à água para que
o líquido que chega à curva já tenha sido concentrado pelo interstício e
cada estágio trabalhe sobre um líquido mais salgado que o anterior. Dois
ramos com a mesma permeabilidade não conseguiriam multiplicar.
\end{solution}

\begin{solution}{exo:b3:renal-osmoregulation:5}
Extremidade aferente: $60 - 18 - 28 = \qty{14}{mmHg}$; extremidade
eferente: $60 - 18 -
35 = \qty{7}{mmHg}$. À medida que a água é filtrada,
as proteínas que ficam se concentram e a pressão oncótica sobe ao longo
do capilar; se ela chegar a \qty{42}{mmHg} a pressão líquida é zero e a
filtração cessa antes do fim (equilíbrio de filtração). A TFG depende
então do fluxo plasmático: um fluxo mais rápido concentra as proteínas
mais devagar, de modo que a filtração continua por um trecho mais longo
e a TFG sobe com o fluxo.
\end{solution}

\begin{solution}{exo:b3:renal-osmoregulation:6}
Carga filtrada: $\qty{40}{mL/min}\times\qty{0.14}{mmol/mL} =
\qty{5.6}{mmol/min}$. Excreção: $70\times 1 = \qty{0.07}{mmol/min}$.
Reabsorvida: $5.53/5.6 = \qty{98.75}{\%}$. Clearance do sódio: $70\times
1/140 = \qty{0.5}{mL/min}$.
\end{solution}

\begin{solution}{exo:b3:renal-osmoregulation:7}
$900/1200 = \qty{0.75}{L}$ por dia. A \qty{400}{mOsm/L}: $900/400 =
\qty{2.25}{L}$ de urina, de modo que cerca de $2.25 + 0.9 = \qty{3.15}{L}$ precisam ser bebidos (menos a água dos alimentos): o
paciente com um mecanismo de concentração falho precisa beber, e corre
perigo se não puder.
\end{solution}

\begin{solution}{exo:b3:renal-osmoregulation:8}
(a) Diabetes insípido: \qtyrange{10}{15}{L} de urina a
\qtyrange{50}{100}{mOsm/L}, sódio plasmático alto sempre que a ingestão
atrasa, sede implacável. (b) Vasopressina inapropriada: urina escassa e
concentrada, água retida, sódio plasmático baixo por diluição, sede não
aumentada embora se continue a beber --- confusão e convulsões quando o
sódio cai muito. (c) Excesso de \omterm{def:b3:renal-osmoregulation:controls}{aldosterona}: sódio retido com água,
volume expandido, pressão arterial alta, potássio perdido (K$^{+}$
plasmático baixo), sódio plasmático perto do normal, já que a água segue
o sal e o rim escapa em parte; volume urinário normal. (d) Dieta sem
sal: a renina e a \omterm{def:b3:renal-osmoregulation:controls}{aldosterona} sobem em um dia, o sódio urinário cai a
alguns milimoles por dia, o sódio plasmático é mantido, e o volume e a
pressão ficam um pouco mais baixos.
\end{solution}

\begin{solution}{exo:b3:renal-osmoregulation:9}
Urina diluída: clearance osmolar $100\times 6/300 = \qty{2}{mL/min}$;
\omterm{prop:b3:renal-osmoregulation:water}{clearance de água livre} $6 - 2 = +\qty{4}{mL/min}$: o rim está
eliminando água sem soluto, como depois de uma sobrecarga hídrica com a
vasopressina desligada. Urina concentrada: clearance osmolar $1000\times 0.6/300 =
\qty{2}{mL/min}$; \omterm{prop:b3:renal-osmoregulation:water}{clearance de água livre} $0.6 - 2 = -\qty{1.4}{mL/min}$: \qty{1.4}{mL/min} de água livre estão sendo
devolvidos ao corpo --- desidratação, vasopressina alta.
\end{solution}

\begin{solution}{exo:b3:renal-osmoregulation:10}
Numere os estágios a partir do córtex; no estágio $k$, chame o
interstício de $I_{k}$. O líquido descendente em $k$ é igual a $I_{k}$; o
ascendente em $k$ é $I_{k} - 200$; e o líquido ascendente em $k$ é o que
saiu da curva e passou pelos estágios $n, n-1, \dots, k+1$. Em regime
estacionário o líquido que entra a \qty{300}{} sai do topo do ramo
ascendente a $I_{1} - 200$, e o interstício de cada estágio supera o do
estágio acima em no máximo o \omterm{prop:b3:renal-osmoregulation:countercurrent}{efeito único}, de modo que $I_{k} \le 300 + 200k$, com igualdade quando todo estágio trabalha em efeito pleno: a
curva alcança $300 + 200n$. O gradiente é, portanto, limitado pelo
comprimento da alça (o número de estágios) vezes o \omterm{prop:b3:renal-osmoregulation:countercurrent}{efeito único}, e não
pelo poder de bombeamento de célula alguma. $n$ é o comprimento da alça
dividido pela distância ao longo da qual o líquido se equilibra com o
interstício: as longas alças justamedulares, e as alças muito longas dos
roedores do deserto, dão $n$ grande; a capacidade da bomba e a lavagem
do soluto pelos vasa recta limitam o efeito.
\end{solution}

\begin{solution}{exo:b3:renal-osmoregulation:11}
\omterm{def:b3:renal-osmoregulation:comparative}{Água metabólica}: $5\times 0.8 = \qty{4}{g}$ de amido dão $4\times 0.6
= \qty{2.4}{g}$. Perdas: urina, $3/5500\ \unit{L} = \qty{0.55}{mL}$;
evaporação, \qty{1.5}{g} (já descontada a recuperação do focinho). Perda
total de \qty{2.05}{g} contra um ganho de \qty{2.4}{g}: uma sobra de
\qty{0.35}{g}, que cobre a pequena perda fecal, e a água higroscópica
das sementes é um bônus. Ele sobrevive sem beber. Sem a recuperação
nasal, a perda respiratória seria de $1.5/0.3 = \qty{5}{g}$, e ele
morreria em poucos dias.
\end{solution}

\begin{solution}{exo:b3:renal-osmoregulation:12}
No fluxo paralelo o sangue e o dialisato se equilibram na metade do
caminho ao longo da membrana e a diferença de concentração que move a
difusão cai a zero: na melhor das hipóteses o sangue sai com metade de
sua ureia. No fluxo em contracorrente o dialisato mais fresco encontra o
sangue mais limpo e o gradiente é mantido ao longo de toda a membrana,
de modo que o sangue pode sair com quase a concentração de entrada do
dialisato --- o mesmo princípio da \omterm{def:b3:renal-osmoregulation:nephron}{alça de Henle}. Clearance: $300\times 0.7 =
\qty{210}{mL/min}$ durante $3\times 4 = \qty{12}{h}$ por semana,
isto é, $210\times 720 = \qty{151}{L}$ por semana; dois rins a
\qty{125}{mL/min} depuram $125\times\num{10080} = \qty{1260}{L}$. A
máquina faz cerca de um oitavo do trabalho dos rins, uma média temporal
de \qty{15}{mL/min} --- o bastante para viver, não o bastante para estar
bem.
\end{solution}

\begin{solution}{pb:b3:renal-osmoregulation:1}
\textbf{1.} Pressão líquida $55 - 15 - 30 = \qty{10}{mmHg}$; TFG $= 12.5
\times 10 = \qty{125}{mL/min}$.
\textbf{2.} $125\times 1440 = \qty{180}{L/d}$; $180/3 = 60$ vezes.
\textbf{3.} Inulina: $62.5\times 1/0.5 = \qty{125}{mL/min}$, igual à
TFG. Creatinina: $1.3\times 1/0.01 = \qty{130}{mL/min}$, um pouco maior
porque o \omterm{def:b3:renal-osmoregulation:nephron}{túbulo proximal} secreta uma pequena quantidade de creatinina
além do que é filtrado.
\textbf{4.} Clearance do PAH $13\times 1/0.02 = \qty{650}{mL/min}$, o
fluxo plasmático renal; fluxo sanguíneo $650/(1 - 0.45) = \qty{1180}{mL/min}$, cerca de \qty{1.2}{L/min}; fração de filtração
$125/650 = 0.19$.
\textbf{5.} Pressão líquida \qty{15}{mmHg}, TFG \qty{188}{mL/min} (menos
na realidade, já que a pressão oncótica sobe ao longo do capilar). A
angiotensina II contrai a arteríola eferente e eleva a pressão
glomerular; bloqueá-la baixa um pouco a pressão e a TFG. No diabetes os
\omterm{def:b3:renal-osmoregulation:nephron}{glomérulos} filtram sob pressão alta e o filtro se desgasta; baixar a
pressão retarda o dano em anos.
\textbf{6.} Os podócitos e seus diafragmas em fenda (ou a carga da
membrana basal): a camada que retém a albumina. Perder \qty{3}{g/d}
baixa a pressão oncótica plasmática; o líquido deixa os capilares em
toda parte e os tecidos incham --- o edema da síndrome nefrótica.
\textbf{7.} Filtrada $125\times 5 = \qty{0.625}{mmol/min}$; excretada
$250\times 1 = \qty{0.25}{mmol/min}$; reabsorvida $0.375$, isto é,
\qty{60}{\%}; clearance $250/5 = \qty{50}{mL/min}$.
\textbf{8.} Filtrados $180\times 0.14 = \qty{25.2}{mol/d}$ de
Na$^{+}$, \qty{580}{g} de sódio, \qty{1.47}{kg} como NaCl. Excretados
$100\times 1.44
= \qty{144}{mmol/d}$; reabsorvidos $1 - 144/\num{25200} = \qty{99.4}{\%}$. Quase um quilo e meio de sal se
perderia por dia.
\textbf{9.} Filtrados $125\times 5 = \qty{0.625}{mmol/min}$, isto é,
\qty{112}{mg/min}, \qty{162}{g/d}; nada é excretado, já que isso está
abaixo do $T_{m}$. A carga filtrada alcança $T_{m}$ em $2.1/0.125 =
\qty{16.8}{mmol/L}$ (\qty{300}{mg/dL}) --- acima do limiar observado, de
cerca de \qty{10}{mmol/L}, porque os \omterm{def:b3:renal-osmoregulation:nephron}{néfrons} diferem e os mais fracos
transbordam primeiro (a dispersão da curva de titulação).
\textbf{10.} Filtrados $125\times 20 = \qty{2.5}{mmol/min}$; excretados
$2.5 - 2.1 = \qty{0.4}{mmol/min}$, \qty{72}{mg/min}, $576$ mmol ou
\qty{104}{g} por dia. Urina extra $576/300 = \qty{1.9}{L}$. A glicose
arrasta água para a urina (diurese osmótica), o paciente se desidrata e
tem sede; \qty{104}{g} de glicose são \qty{1.8}{MJ} perdidos por dia, e
as células, incapazes de captar glicose, queimam gordura e proteína: o
peso cai.
\textbf{11.} Reabsorvidos $\approx \qty{25}{mol/d}$; ATP $25/3 =
\qty{8.4}{mol/d}$; $8.4\times 50 = \qty{420}{kJ}$, cerca de
\qty{6}{\%} de \qty{7}{MJ}.
\textbf{12.} Para secretar cada resíduo o túbulo precisaria de um
transportador que o reconhecesse --- milhares deles, e nenhum para uma
molécula nunca antes encontrada. Filtrar tudo que é pequeno e reabsorver
as poucas centenas de espécies que o corpo quer exige transportadores
apenas para o que se sabe ser valioso; o que não é reconhecido sai por
padrão. O preço é a energia da questão~11.
\textbf{13.} Clearance osmolar $600\times 1/290 = \qty{2.07}{mL/min}$;
\omterm{prop:b3:renal-osmoregulation:water}{clearance de água livre} $1 - 2.07 = -\qty{1.07}{mL/min}$: o rim está
conservando água.
\textbf{14.} Mínimo $600/1200 = \qty{0.5}{L}$; máximo $600/50 =
\qty{12}{L}$; necessidade diária no mínimo $0.5 + 0.9 = \qty{1.4}{L}$.
\textbf{15.} \qty{1000}{mOsm} a \qty{1200}{mOsm/L} exigem \qty{0.83}{L}
de urina: ganho líquido de \qty{0.17}{L} antes do soluto do dia. Depois:
o soluto total do dia é \qty{1600}{mOsm}, \qty{1.33}{L} de urina, mais
\qty{0.9}{L} insensíveis, \qty{2.23}{L} para fora contra \qty{1}{L} para
dentro --- um déficit de \qty{1.23}{L}, contra \qty{1.4}{L} sem beber. O
litro de água do mar melhora o balanço em apenas \qty{0.17}{L}.
\textbf{16.} Urina a \qty{50}{mOsm/L} para o dia: $600/50 =
\qty{12}{L}$; \qty{12.9}{L} precisam ser bebidos a cada dia, cerca de
\qty{13}{L}.
\textbf{17.} $B$ litros de cerveja trazem $20B$ mOsm; a urina máxima é
$(600 + 20B)/50 = 12 + 0.4B$ litros, e ela precisa carregar $B - 0.9$
litros: $B - 0.9 \le 12 + 0.4B$, $B \le \qty{21.5}{L}$. Uns vinte litros
--- mas, se pouco mais for comido, o soluto do dia cai a
\qty{200}{mOsm} e o limite despenca para cerca de \qty{8}{L}: a
hiponatremia dos bebedores pesados que não comem.
\textbf{18.} $298/290 - 1 = \qty{2.8}{\%}$: bem além do \qty{1}{\%} que
dispara a vasopressina e a sede; ambas estão no máximo. Com soluto
constante: $290\times 42 = 298\times V$, $V = \qty{40.9}{L}$, um déficit
de cerca de \qty{1.1}{L}.
\textbf{19.} Seis litros de água diluem o plasma enquanto o suor retira
sal; e o exercício, a dor e o estresse liberam vasopressina
independentemente da \omterm{def:b3:renal-osmoregulation:problem}{osmolaridade}, de modo que o rim não consegue
diluir a urina o bastante para excretar o excesso. O sódio plasmático
cai, as células do encéfalo incham --- convulsões e morte já se
seguiram. O rim precisaria de vasopressina desligada e de \qty{6}{L} de
urina em $U_{\min}$; o corredor deveria beber conforme a sede e repor o
sal.
\textbf{20.} $3/5500\ \unit{L} = \qty{0.55}{mL}$ por dia. Um rim humano
precisaria de $3/1200 = \qty{2.5}{mL}$ para o mesmo soluto: o
rato-canguru usa um quinto da água.
\textbf{21.} Entra: \qty{2.4}{g}. Sai: $1.6 + 0.3 + 0.55 = \qty{2.45}{g}$. Equilibrado com margem de \qty{0.05}{g} por dia, que a
água higroscópica das sementes (alguns por cento de \qty{5}{g})
fornece.
\textbf{22.} Água $500\times 0.65 = \qty{325}{L}$; um quarto é
\qty{81}{L}; a \qty{5}{L/d}, dezesseis dias.
\textbf{23.} $150/800 = \qty{0.19}{L}$, cerca de \qty{190}{mL} de
salmoura --- menos que os \qty{100}{mL} de água do mar mais a água do
peixe, de modo que a ave ganha água. Seu rim, a \qty{600}{mOsm/L}
(\qty{300}{mmol/L} de NaCl), precisaria de \qty{0.5}{L} de urina para
carregar o mesmo sal, mais água do que ela ingeriu: um rim que não vence
a água do mar não pode deixar uma ave bebê-la, e a glândula, ao dobro da
concentração do mar, pode.
\textbf{24.} \qty{30}{g} de água entram, de modo que cerca de
\qty{30}{mL} de urina por dia, a uma pequena fração da \omterm{def:b3:renal-osmoregulation:problem}{osmolaridade} do
sangue (algumas dezenas de mOsm/L contra \num{300}). O sal sai nessa
urina e por difusão pelas brânquias; sem a captação ativa pelas \omterm{def:b3:renal-osmoregulation:comparative}{células
de cloreto} a partir de uma água com um milimol por litro, o peixe se
esgotaria em dias.
\textbf{25.} TFG de \qty{125}{mL/min}, \qty{180}{L} por dia; urina
obrigatória de \qty{0.5}{L} por dia; um litro de água do mar rende
\qty{0.17}{L} líquidos depois de excretado seu sal --- quase nada.
\end{solution}
